\chapter{Ischaemic Heart Disease}
\index{Ischaemic Heart Disease}

\section*{Definition and Overview}
Ischaemic heart disease, also called coronary artery disease or coronary heart disease, is a condition in which the coronary arteries, the vessels that supply blood to the heart muscle, become narrowed or blocked, usually by the build-up of fatty plaques called atherosclerosis. When blood flow is restricted, the heart muscle is deprived of oxygen, causing angina, chest pain or discomfort on exertion, and when a plaque ruptures and a clot blocks an artery completely, a heart attack occurs. Ischaemic heart disease is the leading single cause of death in many countries and worldwide. It is largely preventable and highly treatable through lifestyle change, medicines, and, where needed, procedures to restore blood flow.

\section*{Epidemiology}
Coronary heart disease affects around 580,000 people and causes around 17,000 deaths each year. It is more common in men, and its risk rises steeply with age, although it can occur in younger people, especially with smoking, diabetes, or a strong family history. The major risk factors are smoking, high blood pressure, high cholesterol, diabetes, physical inactivity, obesity, and an unhealthy diet, and Aboriginal and Torres Strait Islander peoples and people from lower socioeconomic backgrounds are disproportionately affected. Rates of heart attack have fallen over recent decades with better prevention and treatment, but the burden remains large.

\section*{Aetiology and Pathophysiology}
Atherosclerosis begins when cholesterol and inflammatory cells accumulate in the walls of the coronary arteries, forming fatty plaques that narrow the vessel. Over years, plaques grow and can become unstable, with a thin cap over a pool of lipids. When such a plaque ruptures, the body forms a blood clot at the site, which can block the artery and cause a heart attack, starving a segment of heart muscle of oxygen. Angina occurs when narrowing limits blood flow during exertion or stress, when the heart needs more oxygen, and resolves with rest. Risk factors damage the artery lining and accelerate this process: smoking, high blood pressure, diabetes, and high cholesterol. The heart muscle that is lost in a heart attack is replaced by scar tissue, impairing the heart's pumping function.

\section*{Clinical Features}
\begin{itemize}
    \item Angina: chest pain or pressure, heaviness, or tightness, often brought on by exertion and relieved by rest
    \item Pain that may spread to the arm, jaw, neck, or back
    \item Breathlessness on exertion
    \item A heart attack: sudden, severe chest pain at rest, lasting more than a few minutes, with sweating, nausea, or breathlessness
    \item Atypical symptoms: fatigue, indigestion-like discomfort, and breathlessness, more common in women, older people, and people with diabetes
    \item Some people have no symptoms until a heart attack, called silent ischaemia
    \item Heart failure with breathlessness and swelling if the pumping function is damaged
\end{itemize}

\section*{Differential Diagnosis}
Chest pain has many causes, and distinguishing heart disease from them matters.
\begin{itemize}
    \item \textbf{Stable angina:} predictable pain on exertion
    \item \textbf{Acute coronary syndrome:} unstable angina and heart attack
    \item \textbf{Other cardiac conditions:} pericarditis, aortic dissection, and arrhythmias
    \item \textbf{Reflux and oesophageal spasm:} common mimics of cardiac pain
    \item \textbf{Musculoskeletal pain:} chest wall, rib, and spinal causes
    \item \textbf{Lung conditions:} pulmonary embolism, pneumonia, and pleurisy
    \item \textbf{Anxiety and panic:} chest tightness and palpitations
\end{itemize}

\section*{When to Seek Medical Care}
Anyone with chest pain or discomfort on exertion should be assessed, and people with risk factors should have their cardiovascular risk reviewed regularly. New, worsening, or unstable symptoms require urgent attention, as do symptoms in people with diabetes, who may have atypical presentations.

\textbf{Contact your local emergency services for:}
\begin{itemize}
    \item Chest pain at rest, lasting more than a few minutes, or with sweating, nausea, or breathlessness
    \item Pain spreading to the arm, jaw, neck, or back
    \item Collapse, confusion, or sudden severe breathlessness
    \item A rapid or irregular heartbeat with chest pain
    \item Any suspicion of a heart attack: do not drive yourself to hospital
\end{itemize}

\section*{Diagnosis and Investigations}
\begin{itemize}
    \item \textbf{History:} the pattern, character, and triggers of chest pain
    \item \textbf{Electrocardiogram (ECG):} detects damage, ischaemia, and rhythm problems
    \item \textbf{Blood tests:} troponin, a marker of heart muscle damage, in suspected heart attack
    \item \textbf{Exercise stress testing:} ECG and imaging during exercise
    \item \textbf{Echocardiogram:} assessment of heart function and wall motion
    \item \textbf{Coronary angiography:} X-ray imaging of the coronary arteries, the definitive test
    \item \textbf{CT coronary angiography:} a non-invasive test for coronary narrowing
    \item \textbf{Cardiovascular risk assessment:} blood pressure, cholesterol, glucose, and risk calculators
\end{itemize}

\section*{Management}
\begin{itemize}
    \item \textbf{Lifestyle change:} smoking cessation, healthy diet, weight loss, and regular activity
    \item \textbf{Medicines:} antiplatelet therapy, statins, blood pressure medicines, and beta-blockers
    \item \textbf{Angina medicines:} nitrates and other medicines to relieve and prevent symptoms
    \item \textbf{Percutaneous coronary intervention:} angioplasty with stenting to open narrowed arteries
    \item \textbf{Coronary artery bypass surgery:} grafting for extensive or complex disease
    \item \textbf{Cardiac rehabilitation:} structured exercise and education after a heart event
    \item \textbf{Management of diabetes and other risk factors:} optimising glucose, blood pressure, and cholesterol
    \item \textbf{Emergency treatment:} rapid re-opening of the blocked artery in a heart attack
\end{itemize}

\section*{Complications}
\begin{itemize}
    \item Heart attack with death of heart muscle
    \item Heart failure from impaired pumping function
    \item Arrhythmias, including ventricular fibrillation and sudden cardiac death
    \item Angina that limits activity and quality of life
    \item Repeat heart attacks and progression of disease
    \item Complications of procedures and surgery
    \item Psychological effects, including anxiety and depression after a heart event
\end{itemize}

\section*{Prognosis and Outcomes}
With modern treatment, most people with ischaemic heart disease live many years and maintain a good quality of life. Stopping smoking, controlling blood pressure and cholesterol, and taking prescribed medicines reduce the risk of heart attack by more than half in high-risk people. Emergency treatment of heart attack is highly effective when delivered early, with survival improving steadily. Cardiac rehabilitation and lifestyle change further improve outcomes. The earlier risk factors are controlled, the better the long-term outlook.

\section*{Prevention and Self-Care}
\begin{itemize}
    \item Do not smoke, and seek help to quit
    \item Eat a heart-healthy diet rich in vegetables, fruit, whole grains, and healthy fats
    \item Be physically active for at least 30 minutes on most days
    \item Maintain a healthy weight
    \item Control blood pressure, cholesterol, and diabetes
    \item Limit alcohol
    \item Take prescribed medicines exactly as directed
    \item Know the warning signs of a heart attack and contact local emergency services immediately
    \item Attend cardiac rehabilitation after a heart event
    \item Manage stress and get support for low mood
\end{itemize}

\section*{Sources and Further Reading}
The following reputable sources provide further information for healthcare students and patients seeking reliable, current advice about ischaemic heart disease.
\begin{itemize}
    \item Heart Foundation Australia. \textit{Coronary heart disease information.} https://www.heartfoundation.org.au/
    \item Healthdirect Australia. \textit{Coronary heart disease.} https://www.healthdirect.gov.au/
    \item National Heart Foundation. \textit{Angina and heart attack information.} https://www.heartfoundation.org.au/
    \item Cardiac Society of Australia and New Zealand. \textit{Coronary disease guidelines.} https://www.csanz.edu.au/
    \item NHS. \textit{Coronary heart disease.} https://www.nhs.uk/conditions/coronary-heart-disease/
    \item World Health Organization. \textit{Cardiovascular diseases.} https://www.who.int/
\end{itemize}
